This master's thesis marks the completion of my studies in the two-year Master's Degree Program in Cybernetics and Robotics at the Norwegian University of Science and Technology (NTNU), Department of Engineering Cybernetics.

The work has been both challenging and rewarding, and has given me the opportunity to apply theory from cybernetics, orbital mechanics, simulation, and control to a complex space mission analysis problem. Developing and analyzing the Basilisk-based simulation framework has provided valuable insight into both the technical modelling process and the practical limitations that arise when building long-horizon spacecraft simulations.

I would like to express my sincere gratitude to my supervisor, Jan Tommy Gravdahl, for his guidance, feedback, and support throughout the thesis work. I would also like to thank my co-supervisor, Corrado Chiatante, for valuable discussions and academic input during the project.

A special thanks goes to Oliver Keving Hasler for valuable input related to the implementation of the simulations in Basilisk. His insights were helpful in addressing several practical modelling and implementation challenges.

Finally, I would like to thank my family, friends, and fellow students for their support, encouragement, and discussions throughout the semester.